\chapter{Contexte et problématique}
\label{chap:problematique}

\section{La chaîne de valeur en trois étages}

Le périmètre du stage couvre la production et la distribution d'acide phosphorique, depuis
l'attaque du minerai jusqu'à la livraison aux consommateurs. Cette chaîne comporte trois
étages.

\begin{figure}[H]
\centering
\begin{tikzpicture}[
  node distance=1.1cm,
  etage/.style={rectangle, draw=ocpbleu, thick, fill=grisclair, rounded corners=2pt,
                text width=9.5cm, align=center, minimum height=1.3cm, font=\small},
  fleche/.style={-{Stealth[length=3mm]}, thick, draw=ocpbleu},
  entree/.style={font=\small\itshape, text=grisfonce},
]
\node[entree] (in) {Minerai de phosphate \; + \; acide sulfurique};
\node[etage, below=0.6cm of in] (e1)
  {\textbf{Étage 1 --- Attaque-filtration}\\[2pt]
   6 lignes : 13AB, 13CD, 13XY, 13ZU, 13E, 13F\\
   produit l'acide à \SI{29}{\percent} --- \SI{8300}{\tonne}~\PdeuxOcinq{}/jour};
\node[etage, below=of e1] (e2)
  {\textbf{Étage 2 --- Concentration}\\[2pt]
   5 lignes : 14AB, 14CD, 14XY, 14ZU, 14EXT\\
   produit l'acide à \SI{54}{\percent} --- \SI{6982}{\tonne}~\PdeuxOcinq{}/jour};
\node[etage, below=of e2] (e3)
  {\textbf{Étage 3 --- Purification}\\[2pt]
   décadmiation \quad$\cdot$\quad clarification \quad$\cdot$\quad cocristallisation};
\node[entree, below=0.6cm of e3] (out)
  {Engrais \quad$\cdot$\quad acide marchand \quad$\cdot$\quad acide purifié};

\draw[fleche] (in) -- (e1);
\draw[fleche] (e1) -- (e2);
\draw[fleche] (e2) -- (e3);
\draw[fleche] (e3) -- (out);
\end{tikzpicture}
\caption{Les trois étages de la chaîne de production}
\label{fig:trois-etages}
\end{figure}

\subsection{Pourquoi concentrer}

L'acide sortant de l'attaque titre environ \SI{26}{\percent} de \PdeuxOcinq{}. C'est
insuffisant pour deux raisons : transporter de l'eau coûte cher, et la plupart des
fabrications d'engrais exigent un acide titrant \SI{50}{\percent} et davantage. On évapore
donc l'eau sous vide.

\subsection{Pourquoi purifier}

Les consommateurs n'ont pas les mêmes exigences de qualité.

\begin{table}[H]
\centering
\caption{Les trois procédés de purification et leur justification}
\label{tab:pourquoi-purifier}
\begin{tabular}{@{}p{4.2cm}p{3.2cm}p{6.4cm}@{}}
\toprule
\textbf{Besoin} & \textbf{Traitement} & \textbf{Raison} \\
\midrule
Engrais destinés au marché européen & Décadmiation &
Le règlement européen plafonne la teneur en cadmium des engrais. \\
\addlinespace
Engrais de qualité, acide marchand & Clarification &
Retirer les solides en suspension (gypse, silice) qui encrassent les installations. \\
\addlinespace
Acide de très haute pureté & Cocristallisation &
Abaisser les impuretés à des niveaux très bas. \\
\bottomrule
\end{tabular}
\end{table}

\section{Le réseau de distribution est incomplet}

Une question se pose naturellement : pourquoi n'importe quelle ligne ne peut-elle pas servir
n'importe quel consommateur ?

La réponse est économique. Une conduite d'acide phosphorique exige des matériaux résistants
--- l'acide attaque la plupart des métaux ---, des pompes, des réchauffeurs (l'acide à
\SI{54}{\percent} est visqueux et fige en refroidissant), des racks de tuyauterie et de la
maintenance. On ne construit une liaison que si elle est régulièrement utile.

Il en résulte un réseau \textbf{partiel}, décrit par deux matrices d'incidence : la matrice
des transferts entre lignes d'acide 29, et la matrice des raccordements ligne
$\to$ consommateur. Ces matrices sont des données du problème, et non des variables.

\section{Ce que fait aujourd'hui un planificateur}

Chaque matin, un ingénieur d'exploitation doit arrêter cinq familles de décisions :

\begin{enumerate}
  \item quelle quantité d'acide 29 envoyer en concentration sur chaque ligne ;
  \item quelles lignes décadmient, et à quel niveau ;
  \item quelle quantité clarifier sur chaque ligne, sachant que les décanteurs sont partagés
        entre deux usages concurrents ;
  \item quels transferts entre ateliers lancer pour éviter débordements et ruptures ;
  \item quelle ligne sert quel consommateur, pour chaque type d'acide.
\end{enumerate}

Ces décisions sont prises au jugement et au tableur. Quatre difficultés en découlent.

\begin{description}
  \item[Le couplage.] Décadmier davantage consomme de la capacité de décantation, donc
        réduit la production d'acide clarifié, donc met en péril la satisfaction d'un autre
        client. Aucune des cinq décisions n'est isolable des autres.
  \item[Les boucles de recyclage.] Les boues des traitements et le raffinat d'EMAPHOS
        reviennent en amont. Le bilan matière n'est donc pas une simple cascade : il comporte
        des rétroactions.
  \item[La combinatoire.] Le nombre de combinaisons possibles dépasse ce qu'un raisonnement
        manuel peut explorer.
  \item[La variabilité humaine.] Le plan retenu dépend de la personne présente ce jour-là,
        ce qui rend les décisions difficilement auditables.
\end{description}

\section{Problématique et objectifs}

\begin{encadre}[title={Problématique}]
Comment déterminer chaque jour, de façon automatique, reproductible et démontrablement
optimale, le plan de production et de distribution d'acide phosphorique qui satisfait
l'ensemble des demandes tout en respectant les contraintes physiques du site ?
\end{encadre}

Cette problématique se décline en quatre objectifs.

\begin{enumerate}
  \item \textbf{Comprendre et formaliser} le procédé industriel en un modèle mathématique
        dont chaque contrainte est reliée à un organe physique identifié.
  \item \textbf{Résoudre} ce modèle à l'optimum global, dans un temps compatible avec un
        usage quotidien.
  \item \textbf{Valider} la solution de manière indépendante, afin de garantir qu'elle est
        physiquement réalisable et non seulement mathématiquement optimale.
  \item \textbf{Restituer} les résultats sous une forme directement exploitable par le
        service.
\end{enumerate}

\section{Ce que le modèle ne décide pas}

Il est aussi important de délimiter ce qui est hors périmètre. Le modèle ne décide ni la
production d'acide 29 (imposée par l'amont), ni les heures de marche des échelons (relevant
du planning de maintenance), ni l'affectation des échelons à la cocristallisation
(configuration), ni les capacités des équipements, ni les recettes des engrais, ni la
demande commerciale.

\begin{encadre}
\textbf{Formulation mathématique de cette distinction.} Tout ce qui précède constitue des
\emph{paramètres} : des données connues avant la résolution. Seules les cinq familles de
décisions de la section précédente sont des \emph{variables}. Établir rigoureusement cette
frontière est la première étape de toute modélisation, et elle conditionne la validité de
tout ce qui suit.
\end{encadre}

\section{Démarche adoptée}

Le travail a été organisé en sept phases, dont l'enchaînement suit une logique simple :
\emph{ne jamais coder ce qui n'a pas d'abord été formalisé}.

\begin{table}[H]
\centering
\caption{Phases du travail}
\label{tab:phases}
\begin{tabular}{@{}clp{7.5cm}@{}}
\toprule
\textbf{Phase} & \textbf{Intitulé} & \textbf{Livrable} \\
\midrule
0 & Cadrage & Système de continuité, feuille de route \\
1 & Compréhension du métier & Documentation du procédé, audit des données \\
2 & Formalisation & Modèle mathématique complet \\
3 & Socle logiciel & Constantes, conversions, scénarios, tests \\
4 & Modèle d'optimisation & Implémentation MILP, résolution \\
5 & Validation et analyse & Validateur indépendant, sensibilité \\
6 & Restitution & Rapports Excel et JSON \\
7 & Rapport & Le présent document \\
\bottomrule
\end{tabular}
\end{table}

Ce découpage n'est pas cosmétique : les phases 3 à 6 ont pu être menées en une seule
itération de travail parce que la phase 2 avait posé le modèle contrainte par contrainte.
L'écriture du code s'est alors réduite à une traduction. Ce point est repris en conclusion.
